\ExerciseSection

\begin{enumerate}
\def\labelenumi{\arabic{enumi})}
\item
  Let X be a Lindelöf space (Chapter I, § 9, Exercise 15), let H be a directed set (with respect to the relation \(\geqslant\) ) of continuous real-valued functions on X such that the lower envelope g of the functions of H is continuous. Show that there exists a decreasing sequence of functions \((f_{n})\) belonging to H which converges pointwise to g.
\item
  Let I be a compact interval in R and let \((f_{n})\) be a sequence of monotone real-valued functions defined on I which converges pointwise in I to a continuous function g. Show that g is monotone and that the sequence \((f_{n})\) converges uniformly to g in I.
\end{enumerate}

¶ 3) Let \(\Gamma\) be a simply transitive group of homeomorphisms of \(\mathbf{R}\) endowed with the topology of pointwise convergence. For each \(x \in \mathbf{R}\) let \(s_x\) denote the element of \(\Gamma\) such that \(s_x(0) = x\) . Show that the mapping \(x \to s_x\) is a homeomorphism of \(\mathbf{R}\) onto \(\Gamma\) {[}use the fact that if \(s \in \Gamma\) is such that \(x < s(x)\) for some \(x \in \mathbf{R}\) , then \(y < s(y)\) for all \(y \in \mathbf{R}]\) ; hence show that \(\Gamma\) is isomorphic to \(\mathbf{R}\) as a topological group {[}use Exercise 17 a) of § 3 and Theorem 1 of Chapter V, § 3{]}.

\begin{enumerate}
\def\labelenumi{\arabic{enumi})}
\setcounter{enumi}{3}
\item
  Let X be a compact space and let H be a vector subspace of C(X; R) such that \(|u| \in H\) whenever \(u \in H\) . Then a continuous real-valued function f on X can be uniformly approximated by functions belonging to H if and only if, for each pair x, y of points of X, the function f satisfies every linear relation \(\alpha f(x) = \beta f(y)\) where \((\alpha, \beta)\) are such that \(\alpha\beta \geqslant 0\) and \(\alpha g(x) = \beta g(y)\) for every function \(g \in H\) . {[}Apply Proposition 2 of no. 1, observing that the image of H under the mapping \(u \to (u(x) u(y))\) is the whole plane \(R^{2}\) , or a line of equation \(\alpha X = \beta Y\) with \(\alpha\beta \geqslant 0\) , or else consists of the single point \((0, 0)\) .{]}
\item
  Let X be a compact space and let H be a set of continuous real-valued functions on X. Let R be the equivalence relation `` \(u(x) = u(y)\) for all \(u \in H\)'' . A continuous real-valued function f on X can then be uniformly approximated by polynomials (resp. polynomials with no constant term) in the functions of H if and only if f is constant on each equivalence class mod R {[}resp. if and only if f is constant on each equivalence class mod R and vanishes at all points x such that \(u(x) = 0\) for all \(u \in H\) {]}.
\end{enumerate}

¶ 6) Let X be a completely regular space and let \(\mathcal{C}^{\infty}(\mathbf{X};\mathbf{R})\) denote the normed subalgebra of \(\mathcal{B}(\mathbf{X};\mathbf{R})\) consisting of all bounded continuous real-valued functions on X. In order that every subalgebra of \(\mathcal{C}^{\infty}(\mathbf{X};\mathbf{R})\) , which separates the points of X and contains the constant functions, should be dense in \(\mathcal{C}^{\infty}(\mathbf{X};\mathbf{R})\) , it is necessary and sufficient that X should be compact. {[}Let \(\beta\mathbf{X}\) be the Stone-Čech compactification of X (Chapter IX, § 1, Exercise 7). Show that if \(\beta\mathbf{X}-\mathbf{X}\) is not empty then there are non-dense subalgebras of \(\mathcal{C}^{\infty}(\mathbf{X};\mathbf{R})\) which separate the points of X and contain the constant functions, by observing that \(\mathcal{C}^{\infty}(\mathbf{X};\mathbf{R})\) may be identified with \(\mathcal{C}(\beta\mathbf{X};\mathbf{R}).]\)

\begin{enumerate}
\def\labelenumi{\arabic{enumi})}
\setcounter{enumi}{6}
\tightlist
\item
  Let X be a non-compact, completely regular space.
\end{enumerate}

\begin{enumerate}
\def\labelenumi{\alph{enumi})}
\tightlist
\item
  Show that the following properties are equivalent:
\end{enumerate}

\(\alpha)\) The subalgebra of \(\mathcal{C}^{\infty}(\mathbf{X};\mathbf{R})\) consisting of functions of the form \(c + f\) , where \(c\) is a constant and \(f\) has compact support (Chapter IX, § 4, no. 3), is dense in \(\mathcal{C}^{\infty}(\mathbf{X};\mathbf{R})\) .

\(\beta)\) If \(\beta X\) is the Stone-Čech compactification of \(X\) , then \(\beta X - X\) consists of a single point.

γ) There is only one uniformity on X, compatible with its topology, with respect to which X is precompact.

δ) If A, B are two completely separated closed subsets of X (Chapter IX, § 1, Exercise 11) then one of A, B is compact.

\(\zeta)\) If \(\mathcal{U}\) is any uniformity on \(X\) compatible with its topology and if \(V\) is any entourage of \(\mathcal{U}\) , there is a subset \(L\) of \(X\) whose complement is compact, such that \(L \times L \subset V\) .

θ) There is only one uniformity on X compatible with its topology. {[}Note that α) implies that X is locally compact and therefore open in βX, and hence show that α) → β); to prove that β) → α) use the Weierstrass-Stone theorem. To show that β) → γ), use the fact that the uniformity induced on X by that of βX is the finest uniformity compatible with the topology of X, with respect to which X is precompact. To show that γ) → δ), use Chapter IX, § 1, Exercise 11. To show that δ) → ζ), show first that if V is any entourage of U, there exists z ∈ X such that V(z) is not compact; otherwise X would be complete with respect to U and paracompact (Chapter II, § 4, Exercise 9), therefore normal; show then with the help of δ) that X would be countably compact (Chapter IX, § 2, Exercise 14), and finally use Chapter

II, § 4, Exercise 6. Then apply this result to an entourage V defined by \(f(x, x') \leqslant 1\) , where \(f\) is a continuous pseudometric on X, and deduce from \(\delta\) ) that there is a subset L of X, whose complement is compact, such that \(L \times L \subset V\) . Finally, to establish that \(\zeta) \Longrightarrow \theta\) ), observe that every ultrafilter on X either converges or is finer than the filter of complements of relatively compact subsets of X.{]}

\begin{enumerate}
\def\labelenumi{\alph{enumi})}
\setcounter{enumi}{1}
\item
  Show that a space X which satisfies the equivalent conditions of a) is pseudo-compact (Chapter IX, § 1, Exercise 22). {[}Use Exercise 12 of Chapter IX, § 1, or else observe that if X is not pseudo-compact, there is a sequence \((x_{n})\) of points of X which has no cluster point and a continuous mapping \(f: \mathbf{X} \to [\mathrm{o}, \mathrm{i}]\) such that \(f(x_{2n}) = \mathrm{o}\) and \(f(x_{2n+1}) = \mathrm{i}\) .{]}
\item
  Let Y be a non-compact, completely regular space. Show that the subspace X of \(\beta Y\) which is the complement of a point of \(\beta Y - Y\) satisfies the conditions of a).
\end{enumerate}

\begin{enumerate}
\def\labelenumi{\arabic{enumi})}
\setcounter{enumi}{7}
\tightlist
\item
  Let \((\mathbf{X}_i)_{1\leqslant i\leqslant n}\) be a finite family of compact spaces, let \(\mathbf{X} = \prod_{i=1}^{n}\mathbf{X}_i\) be their product, and for each index \(i\) let \(\mathbf{F}_i\) be a closed subset of \(\mathbf{X}_i\) . Show that every continuous real-valued function on \(\mathbf{X}\) which vanishes on the union of the sets \(\mathbf{F}_i\times \prod_{j\neq i}\mathbf{X}_j\) ( \(1\leqslant i\leqslant n\) ) can be uniformly approximated by finite sums of functions of the form
\end{enumerate}

\[
(x _ {1}, \dots , x _ {n}) \rightarrow u _ {1} (x _ {1}) \dots u _ {n} (x _ {n}),
\]

where \(u_{i}\) is a continuous real-valued function on \(\mathbf{X}_i\) which vanishes on \(\mathbf{F}_i\) ( \(1 \leqslant i \leqslant n\) ).

* 9) Let \((r_n)_{n \geqslant 1}\) be the sequence of distinct rational numbers belonging to the interval \(\mathbf{I} = [\mathbf{o}, \mathbf{i}]\) of \(\mathbf{R}\) , arranged in some order. Define by induction a sequence of closed intervals \(\mathbf{I}_n \subset \mathbf{I}\) , as follows: \(\mathbf{I}_n\) has as its midpoint the point \(r_{k_n}\) with the smallest index not contained in \(\bigcup_{p < n} \mathbf{I}_p\) ; its length is \(\leqslant \mathbf{i}/4^n\) and it meets no interval \(\mathbf{I}_p\) for which \(p < n\) ; \((\mathbf{I}_n)\) is thus a sequence of mutually disjoint closed intervals. On the product space \(\mathbf{I} \times \mathbf{R}\) we define a continuous real-valued function \(u\) as follows: for each integer \(n \geqslant \mathbf{i}\) , the function \(x \to u(x, n)\) is equal to \(\mathbf{i}\) at an interior point of \(\mathbf{I}_n\) , is equal to \(\mathbf{o}\) on the exterior of \(\mathbf{I}_n\) and takes its values in \([\mathbf{o}, \mathbf{i}]\) ; on the other hand, for each \(x \in \mathbf{I}\) , the function \(y \to u(x, y)\) is affine-linear in each of the intervals \([n, n + \mathbf{i}]\) . Show that \(u\) cannot be uniformly approximated in \(\mathbf{I} \times \mathbf{R}\) by linear combinations of functions of the form \(v(x)w(y)\) , where \(v\) is continuous on \(\mathbf{I}\) and \(w\) is continuous and bounded on \(\mathbf{R}\) . {[}In the space \(\mathcal{C}^\infty(\mathbf{R}; \mathbf{R})\) of continuous bounded real-valued functions on \(\mathbf{R}\) , consider the set of partial functions \(y \to u(x, y)\) as \(x\) runs through \(\mathbf{I}\) ; show that there exists an infinite sequence \((u_n)\) of these functions such that \(||u_n|| = \mathbf{I}\) and \(||u_m - u_n|| = \mathbf{I}\) whenever \(m \neq n\) . Deduce that there cannot exist any finite-dimensional vector subspace \(\mathbf{E}\) of \(\mathcal{C}^\infty(\mathbf{R}; \mathbf{R})\) such that the distance of each \(u_n\) from \(\mathbf{E}\) is \(\leqslant 1/4\) ; for otherwise there would exist a sequence \((v_n)\) of points of \(\mathbf{E}\) such that \(||v_n|| = 2\) and \(||v_n - v_m|| \geqslant 1/2\) whenever \(m \neq n\) , contrary to the fact that every finite-dimensional subspace of \(\mathcal{C}^\infty(\mathbf{R}; \mathbf{R})\) is locally compact.{]} *

\begin{enumerate}
\def\labelenumi{\arabic{enumi})}
\setcounter{enumi}{9}
\tightlist
\item
  Use the Weierstrass-Stone theorem to give a new proof of Urysohn's theorem (Chapter IX, § 4, no. 2, Theorem 2) for closed subspaces of a compact space X. {[}If F is a closed subspace of X, consider the set H of restrictions to F of continuous mappings \(\mathbf{X} \to \mathbf{R}\) , and observe that if \(f \in \mathbf{H}\) and \(|f(x)| \leqslant a\) for all \(x \in \mathbf{F}\) , then there exists a function \(g \in \mathbf{H}\) which coincides with \(f\) on F and is such that \(|g(x)| \leqslant a\) in X.{]}
\end{enumerate}

¶ 11) a) Let X be a completely regular space and let Y be a closed subspace of X; the mapping \(\varphi: u \to u|\mathbf{Y}\) of \(\mathcal{C}^{\infty}(\mathbf{X}; \mathbf{R})\) into \(\mathcal{C}^{\infty}(\mathbf{Y}; \mathbf{R})\) is a continuous linear mapping of norm \(\leqslant 1\) . If X is normal, \(\varphi\) is a surjective strict morphism; and if \(\mathbf{X}_{\mathbf{Y}}\) denotes the quotient space of X obtained by identifying all the points of Y ( \(\mathbf{X}_{\mathbf{Y}}\) is normal by Exercise 15 of Chapter IX, § 4), then the kernel of \(\varphi\) can be identified (as a normed space) with the subspace of \(\mathcal{C}^{\infty}(\mathbf{X}_{\mathbf{Y}}; \mathbf{R})\) consisting of functions which vanish at \(\omega\) , the canonical image of Y in \(\mathbf{X}_{\mathbf{Y}}\) .

\begin{enumerate}
\def\labelenumi{\alph{enumi})}
\setcounter{enumi}{1}
\tightlist
\item
  Suppose that X is metrizable and that the frontier of Y in X is of countable type. Show that \(\varphi: u \to u|Y\) is a strict morphism which has a right inverse (Chapter III, § 6, no. 2, Proposition 3). {[}Let \(d\) be a metric compatible with the topology of X; let \((a_n)\) be a sequence of points of Fr (Y), dense in Fr (Y), and let \(\mathbf{V}_{n,m}\) be the set of points \(x \in \mathbf{X}\) such that \(d(x, a_n) < \mathrm{i}/m\) and \(d(x, \mathbf{Y}) > \mathrm{i}/2m\) ; construct a family of continuous mappings \(f_{n,m}: \mathbf{X} \to [\mathrm{o}, \mathrm{i}]\) such that the support of \(f_{n,m}\) is contained in \(\mathbf{V}_{n,m}\) and such that \(\sum_{n,m} f_{n,m}(x) = \mathrm{i}\) for all \(x \in \complement Y\) , the family \((f_{n,m})\) being uniformly summable in some neighbourhood of each point of \(\complement Y\) . Show then that, for every \(v \in \mathcal{C}^\infty(\mathbf{Y}; \mathbf{R})\) , the function \(u\) which coincides with \(v\) on Y and is equal to \(\sum_{n,m} v(a_n) f_{n,m}(x)\) for all \(x \in \complement Y\) belongs to \(\mathcal{C}^\infty(\mathbf{X}; \mathbf{R})\) and is such that \(\varphi(u) = v.\) {]} (*)
\end{enumerate}

\begin{enumerate}
\def\labelenumi{\arabic{enumi})}
\setcounter{enumi}{11}
\tightlist
\item
  \begin{enumerate}
  \def\labelenumii{\alph{enumii})}
  \tightlist
  \item
    Let X be a compact space, and let \(f: \mathbf{X} \to \mathbf{Y}\) be a continuous surjection of X onto a compact space Y. Show that the mapping \(^{a}f: u \to u \circ f\) of \(\mathcal{C}(\mathbf{Y}; \mathbf{R})\) into \(\mathcal{C}(\mathbf{X}; \mathbf{R})\) is a (norm-preserving) isomorphism of \(\mathcal{C}(\mathbf{Y}; \mathbf{R})\) onto a closed subalgebra of \(\mathcal{C}(\mathbf{X}; \mathbf{R})\) containing the identity element.
  \end{enumerate}
\end{enumerate}

\begin{enumerate}
\def\labelenumi{\alph{enumi})}
\setcounter{enumi}{1}
\item
  Conversely, let A be a closed subalgebra of \(\mathcal{C}(\mathbf{X};\mathbf{R})\) containing the identity element. Show that there exists a continuous surjection \(f\) of \(\mathbf{X}\) onto a compact space \(\mathbf{Y}\) such that \(^a f\) is an isomorphism of \(\mathcal{C}(\mathbf{Y};\mathbf{R})\) onto A. {[}If \((u_{\lambda})_{\lambda \in \mathbf{L}}\) is a family which is dense in A, consider the continuous mapping \(x\to (u_{\lambda}(x))\) of \(\mathbf{X}\) into \(\mathbf{R}^{\mathrm{L}}\) and use the Weierstrass-Stone theorem.{]}
\item
  Deduce from b) that, for each sequence \((u_{n})\) of elements of \(\mathcal{C}(\mathbf{X};\mathbf{R})\) , there exists a metrizable compact space Y and a continuous surjective mapping \(f: X \to Y\) such that the \(u_{n}\) belong to the image of \(\mathcal{C}(\mathbf{Y};\mathbf{R})\) under \(^{a}f\) .
\end{enumerate}

¶ 13) Let X be a compact space and let A denote the normed algebra C(X; R). For each subset M of A, let V(M) denote the set of all \(x \in X\) such that \(u(x) = o\) for all \(u \in M\) . For each subset Y of X, let \(\Im(Y)\) denote the set of all \(u \in A\) such that \(u(x) = o\) for all \(x \in Y\) .

\begin{enumerate}
\def\labelenumi{\alph{enumi})}
\tightlist
\item
  V(M) is a closed subspace of X and \(\Im(Y)\) is a closed ideal of A. If a is the ideal of A generated by a subset M of A, then \(\mathrm{V}(\mathrm{M})=\mathrm{V}(\mathfrak{a})\) ; V and \(\Im\) are order-reversing mappings (with respect to the order relation of inclusion), and we have \(V\{o\}=X\) ; \(\mathrm{V}(\mathrm{A})=\varnothing\) ; \(\mathrm{V}(\{u\})=\varnothing\) for every unit u of A;
\end{enumerate}

\[
\mathrm{V} \left(\bigcup_ {\lambda \in \mathrm{L}} \mathrm{M} _ {\lambda}\right) = \mathrm{V} \left(\sum_ {\lambda \in \mathrm{L}} \mathrm{M} _ {\lambda}\right) = \bigcap_ {\lambda \in \mathrm{L}} \mathrm{V} (\mathrm{M} _ {\lambda})
\]

for every family \((\mathbf{M}_{\lambda})_{\lambda \in \mathbf{L}}\) of subsets of A; \(\mathbf{V}(\mathbf{M}. \mathbf{M}') = \mathbf{V}(\mathbf{M}) \cup \mathbf{V}(\mathbf{M}')\) ; \(\Im(\mathbf{X}) = \{\mathrm{o}\}\) ; \(\Im(\emptyset) = \mathbf{A}\) ; \(\Im\left(\bigcup_{\lambda \in \mathbf{L}} \mathbf{Y}_{\lambda}\right) = \bigcap_{\lambda \in \mathbf{L}} \Im(\mathbf{Y}_{\lambda})\) for every family \((\mathbf{Y}_{\lambda})_{\lambda \in \mathbf{L}}\) of subsets of X.

\begin{enumerate}
\def\labelenumi{\alph{enumi})}
\setcounter{enumi}{1}
\item
  Show that if a is any ideal of A, then \(\Im(\mathrm{V}(a)) = \overline{a}\) , and that if Y is any subset of X, then \(\mathrm{V}(\Im(\mathrm{Y})) = \overline{\mathrm{Y}}\) . (For the first of these relations, observe that \(\mathrm{V}(\overline{a}) = \mathrm{V}(a)\) and that we may therefore suppose a closed; note that if x, y are two distinct points of \(\mathrm{X} - \mathrm{V}(a)\) , there exists \(u \in a\) such that \(u(x) \neq u(y)\) , and use Proposition 4 of no. 2. To prove the second relation, use Urysohn's theorem.)
\item
  Deduce from b) that \(x \to \Im(\{x\})\) is a bijection of X onto the set of maximal ideals of A, which are therefore closed. An ideal of A is closed if and only if it is an intersection of maximal ideals. The ring A has no radical.
\item
  Show that the mapping \(e \to \mathrm{V}(\{e\})\) is a bijection of the set of idempotents of A onto the set of subsets of X which are both open and closed.
\end{enumerate}

A point \(x \in \mathbf{X}\) is isolated if and only if the maximal ideal \(\Im(\{x\})\) is principal.

\begin{enumerate}
\def\labelenumi{\arabic{enumi})}
\setcounter{enumi}{13}
\tightlist
\item
  Let \(\mathbf{X}\) be a topological space. For each function \(u \in \mathcal{C}(\mathbf{X}; \mathbf{R})\) such that \(u(x) > 0\) for all \(x \in \mathbf{X}\) , let \(\mathbf{V}_u\) be the set of all \(f \in \mathcal{C}(\mathbf{X}; \mathbf{R})\) such that \(|f| \leqslant u\) .
\end{enumerate}

\begin{enumerate}
\def\labelenumi{\alph{enumi})}
\item
  Show that the sets \(V_{u}\) are neighbourhoods of o in A = C(X; R) with respect to a Hausdorff topology C compatible with the ring structure of A.
\item
  The topology induced by \(\mathcal{C}\) on \(\mathcal{C}^{\infty}(\mathbf{X};\mathbf{R})\) is finer than the topology of uniform convergence in \(\mathbf{X}\) . These two topologies coincide if and only if \(\mathbf{X}\) is pseudo-compact (Chapter IX, § 1, Exercise 21). If \(\mathbf{X}\) is not pseudo-compact, the topology induced by \(\mathcal{C}\) on the set of constant functions is the discrete topology, so that \(\mathcal{C}\) is not compatible with the \(\mathbf{R}\) -vector space structure of \(\mathbf{A}\) .
\item
  Show that A is a Gelfand ring with respect to the topology \(\mathcal{C}\) (Chapter III, § 6, Exercise 11).
\end{enumerate}

¶ 15) Let X be a completely regular space and let βX be its Stone-Cech compactification. Endow the ring A = C(X; R) with the topology C defined in Exercise 14. For each f ∈ A, let V(f) be the closure in βX of the subset \(\overline{f}^{-1}(\mathrm{o})\) of X. For each subset M of A, let V(M) denote ∩f ∈ M V(f). For each subset Y of βX, let ℑ(Y) denote the set of all f ∈ A such that Y ⊂ V(f), and write ℑ(x) in place of ℑ(\{x\}). We have V(f) = ∅ if and only if f is a unit in A.

\begin{enumerate}
\def\labelenumi{\alph{enumi})}
\item
  If \(f, g\) are two elements of \(\mathbf{A}\) such that \(\overline{f}^{-1}(\mathrm{o}) \cap \overline{g}^{-1}(\mathrm{o}) = \emptyset\) , show that \(\mathrm{V}(f) \cap \mathrm{V}(g) = \emptyset\) . {[}Consider the function \(|f| / (|f| + |g|)\) . Deduce that, for each subset \(\mathbf{M}\) of \(\mathbf{A}\) , we have \(\mathrm{V}(\mathbf{M}) = \mathrm{V}(\mathfrak{a})\) , where \(\mathfrak{a}\) is the ideal generated by \(\mathbf{M}\) in \(\mathbf{A}'\) .
\item
  Show that, for each subset Y of \(\beta X\) , \(\Im(Y)\) is an ideal of A {[}use a){]} and that \(\mathrm{V}(\Im(\mathrm{Y}))\) is the closure \(\overline{Y}\) of Y in \(\beta X\) .
\item
  Let \(u \in \mathbf{A}\) be such that \(u(x) > 0\) for all \(x \in \mathbf{X}\) , and let \(g \in \mathbf{A}\) . Show that there exists \(f \in \mathbf{A}\) such that \(|f - g| \leqslant u\) and such that \(\mathbf{V}(f)\) is a neighbourhood of \(\mathbf{V}(g)\) . {[}Define \(f\) to be equal to \(g + u\) if \(g(x) + u(x) < 0\) , to \(g - u\) if \(g(x) - u(x) > 0\) , and to \(o\) otherwise; consider the function \(f' = \inf ((u + g)^+, (u - g)^-)\) and use \(a)\) .{]}
\item
  Let \(\mathfrak{a}\) be an ideal of \(\mathbf{A}\) . Show that if \(f\in\mathbf{A}\) is such that \(\mathbf{V}(f)\) is a neighbourhood of \(\mathbf{V}(\mathfrak{a})\) , then \(f\in\mathfrak{a}\) . {[}Note first that there exists \(g \in \mathfrak{a}\) such that \(\mathbf{V}(g)\) is contained in the interior of \(\mathbf{V}(f)\) , and then consider the function \(h\) defined on \(\mathbf{X}\) which is equal to \(f(x) / g(x)\) if \(f(x) \neq 0\) , and equal to \(0\) if \(f(x) = 0\) .{]}
\item
  Deduce from c) and d) that \(\Im(\mathrm{V}(\mathfrak{a})) = \overline{\mathfrak{a}}\) for every ideal \(\mathfrak{a}\) of A. {[}Show first that \(\mathrm{V}(\mathbf{M}) = \mathrm{V}(\overline{\mathbf{M}})\) for all subsets M of A, arguing by contradiction.{]} Deduce that, for each subset Y of \(\beta\mathrm{X}\) , \(\Im(\mathrm{Y})\) is a closed ideal of A. {[}Note that \(\Im(\mathrm{Y}) = \Im(\mathrm{V}(\Im(\mathrm{Y})))\) .{]}
\item
  Use e) to show that \(x \to \Im(x)\) is a bijection of \(\beta X\) onto the set of maximal ideals (necessarily closed) of A. An ideal of A is closed if and only if it is an intersection of maximal ideals. The ring A has no radical.
\end{enumerate}

¶ 16) We retain the hypotheses and notation of Exercise 15.

\begin{enumerate}
\def\labelenumi{\alph{enumi})}
\item
  Let a be a closed ideal in the ring A. Show that in the quotient ring A/a the set P of canonical images of functions \(f \geqslant 0\) of A is the set of elements \(\geqslant 0\) with respect to an ordering which makes A/a a lattice-ordered ring {[}observe, using Exercise 15 e), that if f and g belong to a, then so do \(|f| + |g|\) and every function \(h \in A\) such that \(|h| \leqslant |f|\) . The canonical image in A/a of the set of constant functions is isomorphic to R as an ordered field.
\item
  In particular, for each \(x \in \beta X\) , the field \(\mathrm{A}/\Im(x)\) is canonically endowed with an ordered field structure. The maximal ideal \(\Im(x)\) is said to be real if \(\mathrm{A}/\Im(x)\) is isomorphic to R, hyper-real otherwise. For \(\Im(x)\) to be hyper-real it is necessary and sufficient that \(x \in \beta X - X\) and that there is a unit \(f \in A\) such that \(\lim_{y \to x, y \in X} f(y) = 0\) . Deduce that all the maximal ideals of A are real if and only if X is pseudo-compact (Chapter IX, §1, Exercise 21).
\item
  For each \(x \in \beta\mathbf{X}\) , let \(\varphi_x\) be the canonical homomorphism \(\mathbf{A} \to \mathbf{A}/\Im(x)\) . For \(\varphi_x(f)\) to be \(\geqslant 0\) in \(\mathbf{A}/\Im(x)\) , it is necessary and sufficient that there should exist \(g \in \Im(x)\) such that \(f(y) \geqslant 0\) in \(\overline{g}^{-1}(0)\) {[}observe that the relation \(\varphi_x(f) \geqslant 0\) is equivalent to \(f \equiv |f|\pmod{\Im(x)}\) {]}. For \(\varphi_x(f)\) to be \(>0\) , it is necessary and sufficient that there should exist \(g \in \Im(x)\) such that \(f(y) > 0\) in \(\overline{g}^{-1}(0)\) . {[}Note that, if \(f \notin \Im(x)\) , there exists \(g \in \Im(x)\) such that \(\overline{f}^{-1}(0) \cap \overline{g}^{-1}(0) = \emptyset\) , by using Exercise 15.{]}
\item
  Show that if \(\Im(x)\) is hyper-real, the transcendence degree of the field \(\mathbf{A}/\Im(x)\) over \(\mathbf{R}\) is at least equal to \(\mathfrak{c}=\operatorname{Card}(\mathbf{R})\) . {[}Let \(f>0\) be a unit of \(\mathbf{A}\) such that \(\varphi_x(f)=u\) is infinitely large with respect to \(\mathbf{R}\) . Show that as \(r\) runs through the elements of a base of \(\mathbf{R}\) over \(\mathbf{Q}\) (consisting of numbers \(>0\) ), the elements \(\varphi_x(f^r)\) of \(\mathbf{A}/\Im(x)\) are algebraically independent.{]}
\end{enumerate}

* e) For each point \(a = (a_{1}, \ldots, a_{n}) \in \mathbb{R}^{n}\) , let \(f_{a}(\mathbf{X})\) denote the polynomial \(\mathbf{X}^{n} + a_{1}\mathbf{X}^{n-1} + \cdots + a_{n}\) . For each real number \(\xi\) , let \(\nu(\xi)\) be the sum of the multiplicities of the zeros \(z\) of \(f_{a}\) in \(\mathbf{C}\) such that \(\mathrm{R}(z) = \xi\) ; and for each integer \(k\) such that \(\mathrm{i} \leqslant k \leqslant n\) , let \(\rho_{k}(a)\) be the smallest real number \(\xi\) such that \(\sum_{\eta \leqslant \xi} \nu(\eta) \geqslant k\) . Show that each of the functions \(\rho_{k}\) is continuous on \(\mathbb{R}^{n}\) (use Rouché's theorem).

\begin{enumerate}
\def\labelenumi{\alph{enumi})}
\setcounter{enumi}{5}
\item
  Show that, for each \(x \in \beta X\) , the ordered field \(A/\Im(x)\) is real-closed. {[}Let \(F(X) = X^{n} + f_{1}X^{n-1} + \cdots + f_{n}\) be a polynomial of odd degree with coefficients in A; the functions \(g_{k}(y) = \rho_{k}(f_{1}(y), \ldots, f_{n}(y))\) , where the \(\rho_{k}\) are the functions defined in e), are continuous on X, and for each \(y \in X\) there is an index k such that \(F(g_{k}(y)) = 0\) ; deduce that the product \(F(g_{1}) \cdots F(g_{n})\) belongs to \(\Im(x).]\) *
\item
  Let X be an infinite discrete space and let \(y \to F_y\) be a bijection of X onto the set of finite subsets of X; for each \(z \in X\) let \(M_z\) be the set of all \(y \in X\) such that \(z \in F_y\) ; these sets form a base of a filter F on X. Let x be a point of \(\beta X\) such that the corresponding ultrafilter on X (Chapter I, §9, Exercise 27) is finer than F. Let B be a subset of A such that Card (B) ≤ Card (X), and let \(y \to g_y\) be a surjective mapping \(X \to B\) . For each \(z \in X\) put \(f(z) = \mathrm{i} + \sup_{y \in \mathbb{F}_z} g_y(z)\) . Show that \(f(z) > g_y(z)\) for all \(z \in M_y\) , and deduce that \(\varphi_x(f) > \varphi_x(g_y)\) for all \(y \in Y\) {[}use c){]}. Conclude that Card \(A/\Im(x)\) \textgreater{} Card X.
\end{enumerate}

¶ 17) We retain the hypotheses and notation of Exercise 15.

\begin{enumerate}
\def\labelenumi{\alph{enumi})}
\item
  A maximal ideal \(\Im(x)\) is real if and only if, for every infinite sequence \((g_n)\) of elements of \(\Im(x)\) , the intersection of the sets \(\overline{g_n}^{-1}(0)\) is not empty. {[}Use Exercise 16 b); to show that the condition is necessary show that, if \(\bigcap_{n}\overline{g_{n}}^{-1}(0) = \emptyset\) , the function \(f(y) = \sum_{n = 0}^{\infty}\inf (g_n(y),2^{-n})\) is continuous, is a unit in A and tends to o as \(y\) tends to \(x\) while remaining in X; note that \(f(y)\leqslant 2^{-m}\) for \(y\in \bigcap_{k\leqslant m}\overline{g_k}^{-1}(0)\) , and use Exercise 15 a) and the fact that the V(g), where \(g\in \Im(x)\) , form a fundamental system of neighbourhoods of \(x\) in \(\beta \mathbf{X}.\) {]}
\item
  X is said to be real-compact if the only real maximal ideals \(\Im(x)\) are those for which \(x \in \mathbf{X}\) . Show that every completely regular Lindelöf space (Chapter I, § 9, Exercise 15) is real-compact {[}use \(a\) {]}. A completely regular pseudo-compact space is not real-compact unless it is compact.
\item
  Let \(\nu\mathbf{X}\) be the set of all \(x\in \beta \mathbf{X}\) such that \(\Im (x)\) is real. Show that every continuous real-valued function \(f\in \mathbf{A}\) can be extended by continuity to \(\nu\mathbf{X}\) , so that \(\mathcal{C}(\mathbf{X};\mathbf{R})\) and \(\mathcal{C}(\upsilon\mathbf{X};\mathbf{R})\) can be canonically identified. Show that the subspace \(\nu\mathbf{X}\) of \(\beta\mathbf{X}\) is real-compact {[}note that if \(f\) is a unit in \(\mathcal{C}(\mathbf{X};\mathbf{R})\) , then its continuous extension to \(\nu\mathbf{X}\) is a unit in \(\mathcal{C}(\upsilon\mathbf{X};\mathbf{R})\) and that \(\beta(\upsilon\mathbf{X}) = \beta\mathbf{X}]\) . \(\nu\mathbf{X}\) is called the real-compactification of \(\mathbf{X}\) .
\item
  Show that \(\nu\mathbf{X}\) is the completion of \(\mathbf{X}\) with respect to the coarsest uniformity on \(\mathbf{X}\) for which all the functions \(f\in \mathbf{A}\) are uniformly continuous. (Note that a Cauchy filter \(\mathfrak{F}\) on \(\mathbf{X}\) with respect to this uniformity converges to a point \(x\in \beta \mathbf{X}\) ; if \(x\notin \mathfrak{vX}\) , there would be a function \(f\in \mathbf{A}\) which was unbounded in a set of F.{]}
\item
  Deduce from d) that if \(u: X \to Y\) is any continuous mapping of a completely regular space X into a real-compact space Y, then u can be extended by continuity to a mapping \(\upsilon X \to Y\) . The space \(\nu X\) is therefore the solution of the universal mapping problem in which \(\Sigma\) is the structure of a real-compact space, and the \(\alpha\) -mappings and morphisms are continuous mappings (Set Theory, Chapter IV, § 3, no. 1).
\end{enumerate}

\(f\) ) Show that every closed subspace of a real-compact space is real-compact, that every product of real-compact spaces is real-compact, and that in a Hausdorff space any intersection of real-compact subspaces is real-compact. {[}For example, to show that a closed subspace \(\mathbf{X}\) of a real-compact space \(\mathbf{Y}\) is real-compact, consider the extension \(\vec{j}\) to \(\nu\mathbf{X}\) of the canonical injection \(j:\mathbf{X}\to\mathbf{Y}\) , and note that \(j^{-1}(\mathbf{X})=\mathbf{X}\) ; deduce that \(\mathbf{X}\) is closed in \(\nu\mathbf{X}\) . The proofs are similar in the other two cases.{]}

\begin{enumerate}
\def\labelenumi{\alph{enumi})}
\setcounter{enumi}{6}
\tightlist
\item
  Deduce from b), d) and f) that a completely regular space is real-compact if and only if it is homeomorphic to a closed subspace of a product \(R^{I}\) .
\end{enumerate}

\begin{enumerate}
\def\labelenumi{\arabic{enumi})}
\setcounter{enumi}{17}
\tightlist
\item
  Let X be a real-compact space.
\end{enumerate}

\begin{enumerate}
\def\labelenumi{\alph{enumi})}
\item
  Show that, for each non-zero homomorphism \(u: \mathcal{C}(\mathbf{X}; \mathbf{R}) \to \mathbf{R}\) , there exists a unique point \(x \in \mathbf{X}\) such that \(u(f) = f(x)\) for all \(f\) in \(\mathcal{C}(\mathbf{X}; \mathbf{R})\) .
\item
  Let Y be another real-compact space. Show that, for every R-algebra homomorphism \(v: \mathcal{C}(Y; \mathbf{R}) \to \mathcal{C}(X; \mathbf{R})\) which maps identity element to identity element, there exists a unique continuous mapping \(w: X \to Y\) such that \(v(g) = g \circ w\) for all \(g \in \mathcal{C}(Y; \mathbf{R})\) . In particular, \(\mathcal{C}(X; \mathbf{R})\) and \(\mathcal{C}(Y; \mathbf{R})\) are isomorphic if and only if X and Y are homeomorphic.
\end{enumerate}

¶ 19) Let X be a compact space, and let A be the normed C-algebra C(X; C). For each closed ideal a of A, show that the relation \(f \in a\) implies \(\overline{f} \in a\) (note that \(\overline{f}\) is the uniform limit of functions of the form \(g f\) ). Extend the results of Exercise 13 to the algebra A.

¶ 20) a) Let X be a compact space, let A be an R-algebra of finite rank with an identity element, and endow A with the topology defined in Chapter VI, § 1, no. 5. Let B be a subalgebra of the R-algebra C (X; A), and suppose that, for each ε \textgreater{} o, each pair of distinct points x, y of X and each pair of elements u, v of A, there exists f ∈ B such that \textbar\textbar f(x) - u\textbar\textbar{} ≤ ε and \textbar\textbar f(y) - v\textbar\textbar{} ≤ ε (\textbar\textbar{} \textbar\textbar{} being any norm defining the topology of A); and suppose, further, that if u and v belong to R, there is a function f ∈ B satisfying these conditions and taking its values in R. Show that, under these hypotheses, B is dense in C (X; A) with respect to the topology of uniform convergence. {[}Show first that every function of C (X; R) can be uniformly approximated by functions of B and then that the same is true of every constant a ∈ A, by using the hypotheses and a partition of unity.{]}

\begin{enumerate}
\def\labelenumi{\alph{enumi})}
\setcounter{enumi}{1}
\tightlist
\item
  Take A = H, the division ring of quaternions over R. Let B be a subalgebra of the R-algebra \(\mathcal{C}(\mathbf{X}; \mathbf{H})\) such that: (i) for each \(f \in B\) , the conjugate \(\overline{f}\) of f {[}defined by \(\overline{f}(x) = \overline{f(x)}\) for all \(x \in X\) {]} belongs to B; (ii) for each \(\varepsilon > 0\) , each pair of distinct points x, y of X and each pair of elements u, v of H, there exists \(f \in B\) such that \(||f(x) - u|| \leqslant \varepsilon\) and \(||f(y) - v|| \leqslant \varepsilon\) . Show that B is dense in \(\mathcal{C}(\mathbf{X}; \mathbf{H})\) {[}use a){]}. c) Extend the results of Exercise 13 to the algebra \(\mathcal{C}(\mathbf{X}; \mathbf{H})\) , and show in particular that every closed ideal of this algebra is two-sided {[}argue as in Exercise 19, using b){]}.
\end{enumerate}

¶ 21) a) Let X be a totally disconnected compact space, and let A be a topological ring. Endow the ring C (X; A) with the topology of uniform convergence, which is compatible with the ring structure. Let a be a left ideal of C (X; A) and, for each \(x \in X\) , let \(\Im(x)\) be the closure in A of the set of \(f(x)\) for \(f \in a\) ; \(\Im(x)\) is a closed left ideal of A. Show that \(\overline{a}\) is identical with the set of all \(f \in C(X; A)\) such that \(f(x) \in \Im(x)\) for all \(x \in X\) (note that, for each \(x \in X\) , the neighbourhoods of x which are both open and closed form a fundamental system of neighbourhoods of x).

\begin{enumerate}
\def\labelenumi{\alph{enumi})}
\setcounter{enumi}{1}
\tightlist
\item
  Suppose in addition that A has a fundamental system of neighbourhoods of o which are two-sided ideals. Let B be a subring of \(\mathcal{C}(\mathbf{X};\mathbf{A})\) which contains the constants and separates the points of X. Show that B is dense in \(\mathcal{C}(\mathbf{X};\mathbf{A})\) . {[}For each closed subset F of X, each \(x\notin F\) and each neighbourhood U of o in A, show that there exists \(f\in B\) such that \(f(x)=1\) and \(f(y)\in U\) for all \(y\in F\) .{]}
\end{enumerate}

\section*{HISTORICAL NOTE}
\phantomsection
\addcontentsline{toc}{section}{Historical Note}

(Numbers in brackets refer to the bibliography at the end of this note.)

The notion of an arbitrary function was almost unknown at the beginning of the nineteenth century. A fortiori, the idea of studying sets of functions in general and of endowing them with a topological structure did not appear before Riemann's time (see the Historical Note to Chapter I), and it was only towards the end of the nineteenth century that this idea came into systematic use.

Nevertheless, the idea of convergence of a sequence of real-valued functions had been used, more or less consciously, since the beginnings of the infinitesimal calculus. Of course, by convergence we mean here pointwise convergence; other types of convergence could not have been described until the notions of a convergent series and a continuous function had been precisely defined by Bolzano and Cauchy. The latter at first did not recognize the distinction between pointwise convergence and uniform convergence, and believed that he had proved that the sum of any convergent series of continuous functions is continuous ({[}1{]}, (2), vol.~3, p.~120). The error was pointed out almost immediately by Abel, who at the same time showed that every power series is continuous inside its interval of convergence, by a classical argument which, in this particular case, used essentially the notion of uniform convergence ({[}2{]}, pp.~223-224). It remained only to formulate this notion generally, and this was done independently by Stokes and Seidel in 1847-1848, and by Cauchy himself in 1853 ({[}1{]}, (1), vol.~12, p.~30) (*).

Under the influence of Weierstrass and Riemann, the systematic study of the notion of uniform convergence and related questions was developed in the last third of the nineteenth century by the German school (Hankel, du Bois-Raymond) and above all by the Italians; Dini and Arzelà made precise the necessary conditions for the limit of a sequence of continuous functions to be continuous, while Ascoli introduced the fundamental notion of equicontinuity and proved the theorem which characterizes compact sets of continuous functions {[}3{]} (a theorem popularized later by Montel in his theory of ``normal families'', which are relatively compact sets of analytic functions).

On the other hand, Weierstrass himself discovered \([4b]\) the possibility of uniform approximation of a continuous real-valued function in one or more real variables on a bounded set by polynomials. This result immediately aroused lively interest and led to many ``quantitative'' studies (*).

The modern contribution to these questions has been above all to endow them with the full generality of which they are capable, by considering functions whose domain and range are no longer restricted to R or finite-dimensional spaces, and thus placing them in their natural context with the help of general topological concepts. In particular, the theorem of Weierstrass, which had already shown itself to be a powerful weapon in classical analysis, has been extended in recent years by M. H. Stone to much more general situations; developing an idea introduced by H. Lebesgue (in a proof of Weierstrass's theorem) he has shown clearly the important part played in the theory of approximation of continuous real-valued functions by lattices of functions (approximation by ``lattice polynomials'', cf.~§ 4, Proposition 2 and Theorem 2), and has shown also how the generalized Weierstrass theorem has as immediate consequences a whole series of analogous theorems of approximation, which can thus be grouped together in a much more coherent fashion. We have more or less followed his exposition {[}5{]}.

\subsection*{BIBLIOGRAPHY}

{[}1{]} A.-L. CAUCHY, Œuvres, Paris (Gauthier-Villars), 1882-1932.

{[}2{]} N. H. ABEL, Œuvres, vol.~1, p.~223-224, edited by Sylow and Lie (Christiania) 1881.

{[}3{]} G. Ascoli, Sulle curve limiti di una varietà data di curve, Mem. Accad. Lincei (3), vol.~18 (1883), pp.~521-586.

{[}4{]} K. WEIERSTRASS, Mathematische Werke, Berlin (Mayer und Müller), 1894-1903:

\begin{enumerate}
\def\labelenumi{\alph{enumi})}
\item
  Zur Theorie der Potenzreihen, Bd. 1, pp.~67-74;
\item
  Über die analytische Darstellbarkeit sogenannter willkürlicher Functionen reeller Argumente, Bd. 3, pp.~1-37.
\end{enumerate}

{[}5{]} M. H. STONE, The generalized Weierstrass approximation theorem, Mathematics Magazine, 21 (1948), pp.~167-183 and 237-254.

{[}6{]} C. DE LA VALLÉE-POUSSIN, Leçons sur l'approximation des fonctions d'une variable réelle, Paris, (Gauthier-Villars), 1919.

The reference numbers indicate the chapter, section and subsection or exercise, in that order.

T : V, I, 2.

\(a^{x}\) (a any real number \textgreater0, x any real number) : V, 4, I.

\(\log_{a}x\) (a, x real numbers \textgreater0, \(a \neq i\) ): V, 4, I.

\(R^{n}\) : VI, I, I.

\(d(x, y)\) (Euclidean distance) : VI, 2, I.

\(||x||\) (Euclidean norm) : VI, 2, I.

\((x|y)\) (scalar product) : VI, 2, 2.

\(S_{n}, B_{n}, R_{n}^{*}\) : VI, 2, 3.

\(P_{n}(R), P_{n}\) : VI, 3, I.

\(\tilde{R}\) : VI, 3, 3.

\(\infty\) (point of \(\tilde{R}\) ): VI, 3, 3.

\(P_{n,p}(R), P_{n,p}\) : VI, 3, 5.

\(r(G), d(G)\) (G any closed subgroup of \(R^{n}\) ): VII, I, 2.

\(G^{*}\) (G any subgroup of \(R^{n}\) ): VII, I, 3.

\(T^{n}\) : VII, I, 4.

C, i: VIII, I, I.

\(\mathcal{R}(z), \Im(z), \bar{z}, |z|\) (z a complex number): VIII, I, I.

\(C^{*}, U\) : VIII, I, 3.

H : VIII, I, 4.

\(e(x) (= e^{2\pi ix})\) : VIII, 2, I.

\((\widehat{\Delta_{1}, \Delta_{2}})(\Delta_{1}, \Delta_{2}\text{ rays})\) : VIII, 2, 2.

\(\delta, \varpi\) : VIII, 2, 2.

Am (z) (z a complex number): VIII, 2, 2.

cos θ, sin θ, tan θ (θ an angle): VIII, 2, 2.

\(\cos_{a}x\) , \(\sin_{a}x\) , \(\tan_{a}x\) , \(\cot_{a}x\) (x a real number): VIII, 2, 4.

cos x, sin x, tan x, cot x (x a real number): VIII, 2, 4.

\((\widehat{D_{1}, D_{2}})(D_{1}, D_{2}\text{ lines})\) : VIII, 2, 6.

\(\delta_{0}\) : VIII, 2, 6.

\(C^{n}\) : VIII, 4, I.

\(P_{n}(C), \tilde{C}\) : VIII, 4, 3.

∞ (point of \(\tilde{C}\) ): VIII, 4, 3.

\[
\mathbf {P} _ {n, p} (\mathbf {C}): \text { VIII }, 4, 4.
\]

\(\beta \mathbf{X}\) (X a completely regular space): IX, 1, Exercise 7.

\(d(\mathbf{A}, \mathbf{B}), d(x, \mathbf{A})\) {[} \(x\) a point, A, B subsets of a metric space in which the distance between two points is denoted by \(d(x, y)] : \mathrm{IX}, 2, 2\)

\(|x|\) (absolute value in a valued division ring): IX, 3, 2.

\textbar\textbar x\textbar\textbar{} (norm in a normed space) : IX, 3, 3.

D(A) : IX, 5, Exercise 3.

L(C) (C a sieve): IX, 6, 5.

\(\prod_{n \in \mathbb{N}} x_n : \text{IX, Appendix, I.}\)

\(P_{n=h} x_n : \text{IX, Appendix, 3.}\)

\(\prod_{n=0}^{n} x_n: \text{IX, Appendix, 3.}\)

\(\mathcal{F}(\mathrm{X};\mathrm{Y}),\mathrm{H}(x)\) {[}H a subset of \(\mathcal{F}(\mathrm{X};\mathrm{Y})]\) , \(\Phi (x)\) {[} \(\Phi\) a filter on \(\mathcal{F}(\mathrm{X};\mathrm{Y})]\) , \(u|\mathrm{A},\mathrm{H}|\mathrm{A}\) {[}A a subset of \(\mathrm{X},u\in \mathcal{F}(\mathrm{X};\mathrm{Y}),\mathrm{H}\subset \mathcal{F}(\mathrm{X};\mathrm{Y})]:\mathrm{X},\) 1.

\(W(\mathbf{V})\) (V an entourage of \(\mathbf{Y}\) ): \(\mathbf{X}, \mathrm{I}, \mathrm{I}\) .

\[
\mathcal {F} _ {u} (\mathrm{X}; \mathrm{Y}): \mathrm{X}, \mathrm{I}, \mathrm{I}.
\]

\(\mathcal{F}_{\mathfrak{S}}(\mathbf{X};\mathbf{Y}):\mathbf{X},\mathbf{1},\mathbf{2}.\)

\(W(A, V) (A \in \mathfrak{S}, V \text{ an entourage of } Y): X, I, 2.\)

\[
\mathcal {F} _ {s} (\mathrm{X}; \mathrm{Y}), \mathcal {F} _ {c} (\mathrm{X}; \mathrm{Y}): \mathrm{X}, \mathrm{I}, 3.
\]

\[
\mathcal {C} (\mathbf {X}; \mathbf {Y}), \mathcal {C} _ {\mathfrak {S}} (\mathbf {X}; \mathbf {Y}), \mathcal {C} _ {s} (\mathbf {X}; \mathbf {Y}), \mathcal {C} _ {c} (\mathbf {X}; \mathbf {Y}), \mathcal {C} _ {u} (\mathbf {X}; \mathbf {Y}): \mathbf {X}, \mathrm{I}, 6.
\]

\[
\tilde {\mathcal {C}} _ {\mathfrak {S}} (\mathrm{X}; \mathrm{Y}): \mathrm{X}, \mathrm{I}, 6.
\]

\(\tilde{x}\) (x a point of X): X, 2, 1.

\(\mathcal{B}_{\mathfrak{S}}(\mathbf{X};\mathbf{Y}),\mathcal{B}(\mathbf{X};\mathbf{Y})\) (Y a metric space): \(\mathbf{X},3,\mathrm{i}\)

\(||u||\) (u a bounded mapping into a normed space): X, 3, 2.

\(\mathscr{L}\left(\mathbf{X}_{1},\ldots,\mathbf{X}_{n};\mathbf{Y}\right)\left(\mathbf{X}_{1},\ldots,\mathbf{X}_{n}\text{ and Y being normed spaces}\right):\mathbf{X},3,2.\)

\(||u||\) (u a multilinear mapping of a product of normed spaces into a normed space): X, 3, 2.

\(T(\mathbf{K},\mathbf{U})\) (K a compact subset of \(\mathbf{X}\) , U an open subset of \(\mathbf{Y}\) ): \(\mathbf{X}\) , 3, 4. \(\mathcal{C}_{\beta}:\mathbf{X}\) , 3, 5.

\(\mathcal{C}^{\infty}\) (X; R): X, 4, Exercise 6.

\(\nu X:X,4,\) Exercise 17.

The reference numbers indicate the chapter, section and subsection or Exercise, in that order.

Abscissa : VI, 1, 4. Abscissas, axis of : VI, 1, 4. Absolute value : IX, 3, 2. Absolute value, improper : IX, 3, 2. Absolute value, p-adic : IX, 3, 2. Absolute values, equivalent : IX, 3, 2. Absolute value of a complex number : VIII, 1, 1. Absolutely convergent infinite product of complex numbers : VIII, 3, 3. Absolutely convergent series (in a normed space) : IX, 3, 6. Absolutely convergent series of points in R\^{}n : VII, 3, 2. Absolutely summable family (in a normed space) : IX, 3, 6. Acute angular sector : VIII, 2, 5. Additive group of R\^{}n : VI, 1, 2. Additive group of real numbers mod a : V, 1, 2. Additive uniformity of R\^{}n : VII, 1, 2. Admissible subset : IX, 6, Exercise 12. Affine linear varieties, orthogonal : VI, 2, 2. Algebra, normed : IX, 3, 7. Algebraic norm of a complex number : VIII, 1, 1. Algebraic number : VIII, 1, Exercise 2. Almost-open mapping : IX, 5, Exercise 24. Almost-open set : IX, 5, Exercise 6. Almost-periodic function : X, 3, Exercise 28. Almost-periodic (point) : X, 3, Exercise 27. Amplitude of a complex number : VIII, 2, 2 and VIII, 2, 3. Angle, flat : VIII, 2, 2. Angle of a pair of rays : VIII, 2, 2 and VIII, 2, 3. Angle of an angular sector : VIII, 2, 5. Angle, positive right : VIII, 2, 2.

Angular sector (acute, right, obtuse, flat, salient, re-entrant): VIII, 2, 5. Approximation, uniform: X, 4, 1. Archimedean (linearly ordered group): V, 3, Exercise 1. Archimedes' axiom: V, 2. Argument of a complex number: VIII, 2, 2. Ascoli's theorem: X, 2, 5. Associated subgroup of a subgroup of \(\mathbf{R}^n\) : VII, 1, 3. Axiom of Archimedes: V, 2. Axis, coordinate: VI, 1, 4. Axis, imaginary: VIII, 1, 2. Axis of abscissas: VI, 1, 4. Axis of ordinates: VI, 1, 4. Axis, real: VIII, 1, 2.

Compact convergence, uniformity of : X, 1, 3.

Compact-open topology : X, 3, 4.

Compactification, Stone-Čech : IX, 1, Exercise 7.

Compatible (metric and topology) : IX, 2, 5.

Compatible (metric and uniformity) : IX, 2, 4.

Compatible (norm and algebra structure) : IX, 3, 7.

Completely normal space : IX, 4, Exercise 3.

Completely regular filter : IX, 1, Exercise 8.

Completely regular space : IX, 1, 5.

Completely regular space associated with a uniformizable space : IX, 1, Exercise 4.

Completely separated (closed sets in a completely regular space): IX, 1, Exercise 11.

Complex hyperplane in \(\mathbf{C}^n\) : VIII, 4, 1.

Complex linear variety in \(\mathbf{C}^n\) : VIII, 4, 1.

Complex lines in \(\mathbf{C}^n\) : VIII, 4, 1.

Complex number : VIII, 1, 1.

Complex number space of n dimensions : VIII, 4, 1.

Complex planes in \(\mathbf{C}^n\) : VIII, 4, 1.

Complex projective line : VIII, 4, 3.

Complex projective plane : VIII, 4, 3.

Complex projective space of n dimensions : VIII, 4, 3.

Conjugate of a complex number : VIII, 1, 1.

Continuous partition of unity : IX, 4, 3.

Convergent infinite product (in a normed algebra) : IX, Appendix, 3.

Convergent, normally : X, 3, 2.

Convergent, pointwise : X, 1, 3.

Convergent, uniformly : X, 1, 1.

Coordinate axis : VI, 1, 4.

Coordinate variety : VI, 1, 4.

Cosine of an angle : VIII, 2, 2.

Cosine of a number : VIII, 2, 4.

Cotangent of an angle : VIII, 2, 2.

Cotangent of a number : VIII, 2, 4.

Countable type (metrizable space) : IX, 2, 8.

Countably compact space : IX, 2, Exercise 14.

Covering, divisible : IX, 4, Exercise 16.

Covering, even : IX, 4, Exercise 16.

Covering, point-finite : IX, 4, 3.

Cross of a pair of lines : VIII, 2, 6.

Cross, right : VIII, 2, 6.

Cube : IX, 1, 5.

Cube (open, closed) : VI, I, I.

Curve, Peano : VI, 1, Exercise 2.

Degree (unit of angular measure) : VIII, 2, 3.

Diameter : IX, 2, 2.

Diametral hyperplane : VI, 2, 4.

Diametrically opposite points of \(\mathbf{S}_n:\mathbf{VI},3,\mathbf{i}\)

Dimension of a closed subgroup of \(\mathbf{R}^n\) : VII, 1, 2.

Dini's theorem : X, 4, 1.

Direction ratios of a line or ray : VI, 1, 4.

Direction vector of a line or ray : VI, 1, 4.

Disc (open, closed) : VI, 2, 3.

Discrete family of subsets : IX, 4, Exercise 18.

Displacement, Euclidean : VI, 2, 2.

Distance between two sets : IX, 2, 2.

Distance, Euclidean : VI, 2, 1.

Distance from a point to a set : IX, 2, 2.

Divisible covering : IX, 4, Exercise 16.

Division ring of quaternions : VIII, 1, 4.

Division ring, valued : IX, 3, 2.

\(\varepsilon\) -period (of a function): X, 3, Exercise 28.

Equicontinuous set of mappings : X, 2, 1.

Equicontinuous at a point : X, 2, 1.

Equicontinuous, uniformly : X, 2, 1.

Equipartition mod i : VII, i, Exercise 14.

Equivalent absolute values : IX, 3, 2.

Equivalent families of pseudometrics : IX, 1, 2.

Equivalent norms : IX, 3, 3.

Equivalent pseudometrics : IX, 1, 2.

Equivalent semi-absolute values : IX, 3, Exercise 10.

Essential mapping into \(\mathbf{P}_n\) : VI, 3, Exercise 2.

Essential mapping into \(\mathbf{S}_1\) : VI, 2, Exercise 6.

Euclidean ball (closed, open) : VI, 2, 3.

Euclidean displacement : VI, 2, 2.

Euclidean distance : VI, 2, 1.

Euclidean norm in \(\mathbf{R}^n\) : VI, 2, 1.

Euclidean sphere : VI, 2, 3.

Even covering : IX, 4, Exercise 16.

Exponential function : V, 4, 1.

Factors of a product (in a normed algebra) : IX, Appendix, 1.

Family, absolutely summable : IX, 3, 6.

Family of functions subordinate to a family of subsets : IX, 4, 3.

Farey series : VII, 1, Exercise 13.

Field of complex numbers : VIII, 1, 1.

Filter, completely regular : IX, 1, Exercise 8.

Filter, maximal completely regular : IX, 1, Exercise 8.

Finite open coverings, uniformity of : IX, 4, Exercise 17.

Flat angle : VIII, 2, 2.

Flat angular sector : VIII, 2, 5.

Function, almost-periodic : IX, 3, Exercise 28.

Function, exponential : V, 4, 1.

Function, periodic (defined on \(\mathbf{R}^n\) ): VII, 1, 6.

Function, q-ply periodic : VII, 1, 6.

Generated by a set of subsets (σ-algebra) : IX, 6, 3.

Grade (unit of angular measure) : VIII, 2, 3.

Group, additive (of \(\mathbf{R}^n\) ): VI, 1, 2.

Group, additive (of real numbers mod a): V, 1, 2.

Group, metrizable : IX, 3, 1.

Group of Euclidean displacements : VI, 2, 2.

Group, one-parameter : V, 3.

Group, orthogonal : VI, 2, 2.

Half-line : VI, 1, 4.

Half-spaces (open, closed) determined by a hyperplane : VI, 1, 4.

Hausdorff semi-absolute value : IX, 3, Exercise 9.

Hausdorff semi-absolute value associated with a semi-absolute value :

Hemisphere (open, closed) : VI, 2, 4.

Homographic function : VI, 3, 5.

Hyper-real maximal ideal : X, 4, Exercise 16.

Hyperplane at infinity : VI, 3, 3.

Hyperplane, complex (in \(\mathbf{C}^n\) ): VIII, 4, 1.

Hyperplane, diametral : VI, 2, 4.

Hyperplane of projection (in a stereographic projection) : VI, 2, 4.

Imaginary axis : VIII, 1, 2.

Lattice in \(\mathbf{R}^n\) : VII, I, I.

Left-invariant metric : IX, 3, 1.

Left-invariant pseudometric : IX, 3, Exercise 1.

Length of a segment in \(\mathbf{R}^n:\mathrm{VI},2,\mathrm{i}\)

Line, broken : VI, 1, Exercise 6.

Line, simple broken : VI, 1, Exercise 9.

Line, complex projective : VIII, 4, 3.

Line, real projective : VI, 3, 1.

Linear combination of functions with coefficients in a normed space: X, 4, 4.

Linear variety, complex (in \(\mathbf{C}^n\) ): VIII, 4, 1.

Linearly accessible : IX, 6, Exercise 11.

Lines, complex (in \(\mathbf{C}^n\) ): VIII, 4, 1.

Locally paracompact space : IX, 4, Exercise 27.

Logarithm to base \(a:\mathbf{V},4,\mathrm{i}\)

Logarithmic spiral : VIII, 3, Exercise 5.

Lusin space : IX, 6, 4.

Mapping, almost open : IX, 5, Exercise 24.

Mapping, Borel : IX, 6, Exercise 16.

Mapping induced by a sifting : IX, 6, 5.

Mapping into \(\mathbf{P}_n\) , essential or inessential : VI, 3, Exercise 2.

Mapping into \(S_{1}\) , essential or inessential : VI, 2, Exercise 6.

Mapping, isometric : IX, 2, 2.

Mapping, piecewise linear : VI, 1, Exercise 6.

Maximal completely regular filter : X, 1, Exercise 8.

Meagre set : IX, 5, 2.

Measure of an angle : VIII, 2, 3.

Measure of a cross : VIII, 2, 6.

Measure, principal (of an angle): VIII, 2, 3.

Measure, principal (of a cross): VIII, 2, 6.

Metacompact space : IX, 4, Exercise 25.

Metric : IX, 2, 1.

Metric associated with a pseudometric : IX, 2, 1.

Metric compatible with a topology : IX, 2, 5.

Metric compatible with a uniformity : IX, 2, 4.

Metric, left (right) invariant : IX, 3, 1.

Metric space : IX, 2, 1.

Metrizable topological group : IX, 3, 1.

Metrizable topological space : IX, 2, 5.

Metrizable topological space of countable type : IX, 2, 8.

Metrizable uniform space : IX, 2, 4.

Metrizable uniformity: IX, 2, 4.

Multipliable sequence (in a normed algebra) : IX, Appendix, 1.

Nagata-Smirnov theorem : IX, 4, Exercise 22.

Negative real semi-axis : VIII, 1, 2.

Non-meagre space : IX, 5, Exercise 7.

Norm (on a vector space): IX, 3, 3.

Norm, algebraic (of a complex number): VIII, 1, 1.

Norm compatible with an algebra structure : IX, 3, 7.

Norm, Euclidean (in \(\mathbf{R}^n\) ): VI, 2, 1.

Normal space : IX, 4, 1.

Normally convergent series : X, 3, 2.

Normed algebra : IX, 3, 7.

Normed space : IX, 3, 3.

Norms, equivalent : IX, 3, 3.

Nowhere dense set : IX, 5, 1.

Number, algebraic : VIII, 1. Exercise 2.

Number, complex : VIII, 1, 1.

Number space of dimension n, complex : VIII, 4, 1.

Number space of dimension n, real : VI, I, I.

Obtuse angular sector : VI, 2, 5.

One-parameter groups : V, 3.

Open ball : IX, 2, 2.

Opposite rays : VI, 1, 4.

Ordered partition : IX, Appendix, Exercise 2.

Ordinate : VI, I, 4.

Ordinates, axis of : VI, 1, 4.

Origin of a ray : VI, 1, 4.

Orthogonal affine linear varieties : VI, 2, 2.

Orthogonal group : VI, 2, 2.

Orthogonal transformation : VI, 2, 2.

Orthogonal vectors, vector subspaces : VI, 2, 2.

Oscillation of a function at a point : IX, 2, 3.

Oscillation of a function in a set : IX, 2, 3.

\(p\) -adic absolute value: IX, 3, 7.

Parallelotope (open, closed) : VI, 1, 3.

Partition of unity, continuous : IX, 4, 3.

Partition, ordered : IX, Appendix, Exercise 2.

Peano curve : VI, 1, Exercise 2.

Perfectly normal space : IX, 4, Exercise 7.

Periodic function, defined on \(\mathbf{R}^n\) : VII, 1, 6.

Periods of a periodic function : VII, 1, 6.

Piecewise linear mapping : VI, I, Exercise 6.

Pigeon-hole principle : VII, 1, Exercise 11.

Plane, complex projective : VIII, 4, 3.

Plane, real projective : VI, 3, 1. Planes, complex (in C\(^{n}\)) : VIII, 4, 1. Point at infinity : VI, 3, 3. Point-finite covering : IX, 4, 3. Point of uniform convergence : X, 1, Exercise 9. Pointwise convergence in a subset of X : X, 1, 3. Pointwise convergence, topology of : X, 1, 3 and 3, 4. Pointwise convergence, uniformity of : X, 1, 3. Pointwise convergent (filter) : X, 1, 3. Polish space : IX, 6, 1. Polynomial in functions belonging to a given set : X, 4, 2 and 4, 4. Polynomial, trigonometric : X, 4, 4. Positive real semi-axis : VIII, 1, 2. Precompact convergence, uniformity of : X, 1, 3. Principal measure of an angle : VIII, 2, 3. Principal measure of a cross : VIII, 2, 6. Principal system of periods of a q-ply periodic function : VII, 1, 6. Product, absolutely convergent (of complex numbers) : VIII, 3, 3. Product, infinite : IX, Appendix, 3. Product of a multipliable sequence in a normed algebra : IX, Appendix, 1. Product, scalar : VI, 2, 2. Projection, central : VI, 2, 3. Projection, hyperplane of : VI, 2, 4. Projection, stereographic : VI, 2, 4. Projection, vertex of : VI, 2, 4. Projective space of dimension n, complex : VIII, 4, 3. Projective space of dimension n, real : VI, 3, 1. Proper at a point (group of operators) : X, 2, Exercise 18. Pseudo-compact space : IX, 1, Exercise 22. Pseudometric : IX, 1, 1. Pseudometric, invariant : IX, 3, Exercise 1. Pseudometrics, equivalent : IX, 1, 2. Pure imaginary (complex number) : VIII, 1, 1.

q-ply periodic function on R\(^{n}\) : VII, 1, 6. Quadric (in P\(_{n}\)) : VI, 3, Exercise 10. Quadric (in R\(^{n}\)) : VI, 2, Exercise 10. Quadric cone (in P\(_{n}\)) : VI, 3, Exercise 11. Quadric cone (in R\(^{n}\)) : VI, 2, Exercise 11. Quaternion : VIII, 1, 4.

Radian, measure : VIII, 2, 3. Radius of a ball : IX, 2, 2. Radius of a ball or sphere : VI, 2, 3 and IX, 2, 2.

Rational rank of a subgroup of \(\mathbf{R}^n\) : VII, 1.

Ratios, direction : VI, 1, 4.

Ray (closed, open) : VI, 1, 4.

Rays, opposite : VI, 1, 4.

Real axis : VIII, 1, 2.

Real-compact space : X, 4, Exercise 17.

Real compactification of a space : X, 4, Exercise 17.

Real linear variety (in \(\mathbf{C}^n\) ): VIII, 4, i.

Real maximal ideal: X, 4, Exercise 16.

Real number space of \(n\) dimensions : VI, I, I.

Real part of a complex number : VIII, 1, 1.

Real projective line : VI, 3, 1.

Real projective plane : VI, 3, 1.

Real projective space of \(n\) dimensions : VI, 3, 1.

Real valuation of a division ring : IX, 3, 2.

Rectangle (open, closed) : VI, I, I.

Re-entrant angular sector : VIII, 2, 5.

Right angle, positive : VIII, 2, 2.

Right angular sector : VIII, 2, 5.

Right cross : VIII, 2, 6.

Right-invariant metric : IX, 3, 1.

Right-invariant pseudometric : IX, 3, Exercise 1.

\(\sigma\) -algebra: IX, 6, 3.

Salient angular sector : VIII, 2, 5.

Saturated family of pseudometrics : IX, 1, 2.

Saturated set of subsets : X, 1, Exercise 5.

Scalar product : VI, 2, 2.

Sector, angular : VIII, 2, 5.

Segment (closed, open, open at x and closed at y) : VI, 1, 4.

Segment, length of : VI, 2, 1.

Semi-absolute value : IX, 3, Exercise 9.

Semi-absolute value, Hausdorff : IX, 3, Exercise 9.

Semi-absolute values, equivalent : IX, 3, Exercise 10.

Semi-axis, negative real : VIII, 1, 2.

Semi-axis, positive real : VIII, 1, 2.

Semicircle (open, closed) : VI, 2, 4.

Separable metrizable space : IX, 2, 8.

Separating set of functions : X, 4, 1.

Separation of points of a set by a set of functions : X, 4, 1.

Sequence, multipliable : IX, Appendix, 1.

Series, absolutely convergent (of points of \(\mathbf{R}^n\) ): VII, 3, 2.

Series, Farey : VII, 1, Exercise 13.

Set, almost open : IX, 5, Exercise 6.

Set, Borel : IX, 6, 3.

Set, bounded : IX, 2, 2.

Set, bounded (in \(\mathbf{R}^n\) ): VI, I, I.

Set, capacitable : IX, 6, 9.

Set, meagre : IX, 5, 2.

Set, nowhere dense : IX, 5, 1.

Set, Souslin : IX, 6, 2.

Set, starlike (in \(\mathbf{R}^n\) ): VI, 2, Exercise 12.

Set, thin : IX, 5, Exercise 2.

Shell of a starlike set : VI, 2, Exercise 12.

Side of a cube : VI, I, I.

Sides of a broken line : VI, 1, Exercise 6.

Sieve : IX, 6, 5.

Sifting : IX, 6, 5.

Sifting, mapping induced by : IX, 6, 5.

Simple broken line : VI, 1, Exercise 9.

Sine of an angle : VIII, 2, 2.

Sine of a number : VIII, 2, 4.

Slice (in a linearly ordered set): IX, Appendix, Exercise 2.

Slope of a line : VIII, 2, 6.

Souslin set : IX, 6, 2.

Souslin space : IX, 6, 2.

Space associated with a sifting : IX, 6, 5.

Space, Baire : IX, 5, 3.

Space, collectively normal : IX, 4, Exercise 18.

Space, completely normal : IX, 4, Exercise 3.

Space, completely regular : IX, 1, 5.

Space, complex number : VIII, 4, 1.

Space, complex projective : VIII, 4, 3.

Space, countably compact : IX, 2, Exercise 14.

Space, locally paracompact : IX, 4, Exercise 27.

Space, Lusin : IX, 6, 4.

Space, metacompact : IX, 4, Exercise 25.

Space, metric : IX, 2, 1.

Space, metrizable of countable type : IX, 2, 8.

Space, metrizable topological : IX, 2, 5.

Space, metrizable uniform : IX, 2, 4.

Space, non-meagre : IX, 5, Exercise 7.

Space, normal : IX, 4, 1.

Space, normed : IX, 3, 3.

Space of loops : X, 3, 4.

Space of paths : X, 3, 4.

Space of projective linear varieties of dimension \(p\) in complex projective space of dimension \(n: \text{VIII}, 4, 3\) .

Space of projective linear varieties of dimension \(p\) in real projective space of dimension \(n: \mathrm{VI}, 3, 5\) .

Space, perfectly normal : IX, 4, Exercise 7.

Space, Polish : IX, 6, 1.

Space, pseudo-compact : IX, 1, Exercise 21.

Space, real-compact : X, 4, Exercise 17.

Space, real number : VI, I, I.

Space, real projective : VI, 3, 1.

Space, separable metrizable : IX, 2, 8.

Space, Souslin : IX, 6, 2.

Space, strongly zero-dimensional : IX, 6, Exercise 1.

Space, submetrizable : IX, 2, Exercise 23.

Space, totally non-meagre : IX, 5, Exercise 12.

Space, ultrametric : IX, 2, Exercise 4.

Space, zero-dimensional : IX, 6, 4.

Sphere : IX, 2, 2, VI, 2, 3.

Spiral, logarithmic : VIII, 3, Exercise 5.

Square (closed, open) : VI, I, I.

Starlike set in \(\mathbf{R}^n\) : VI, 2, Exercise 12.

Stereographic projection : VI, 2, 4.

Stone-Čech compactification : IX, 1, Exercise 7.

Stone's theorem : X, 4, 1.

Stone-Weierstrass theorem : X, 4, 2.

Strict sifting : IX, 6, 5.

Strongly zero-dimensional space : IX, 6, Exercise 1.

Subgroup, associated : VII, 1, 3.

Submetrizable space : IX, 2, Exercise 23.

Subordinate to a family of subsets (family of functions) : IX, 4, 3.

Support of a function : IX, 4, 3.

Tangent of an angle : VIII, 2, 2.

Tangent of a cross : VIII, 2, 6.

Tangent of a number : VIII, 2, 4.

Theorem, Ascoli's : X, 2, 5.

Theorem, Baire's : IX, 5, 3.

Theorem, Dini's : X, 4, 1.

Theorem of R. Ellis : X, 3, Exercise 25.

Theorem of Nagata-Smirnov : IX, 4, Exercise 22.

Theorem, Stone's : X, 4, 1.

Theorem, Stone-Weierstrass : X, 4, 2.

Theorems, Urysohn's : IX, 4, 1 and 4, 2.

Thin set : IX, 5, Exercise 2.

Topology, compact-open : X, 3, 4.

Topology of compact convergence : X, 1, 3 and 3, 4.

Topology of pointwise convergence : X, 1, 3.

Topology of pointwise convergence in a subset of X : X, 1, 3.

Topology of S-convergence : X, 1, 2.

Topology of uniform convergence : X, 1, 1.

Topology of uniform convergence in a subset of X : X, 1, 3.

Torus, one-dimensional : V, 1, 2.

Torus, n-dimensional : VII, 1, 4.

Torus of revolution : VII, 1, Exercise 15.

Totally non-meagre space : IX, 5, Exercise 12.

Transformation, orthogonal : VI, 2, 2.

Triangle inequality : IX, 1, 1.

Triangle inequality : VI, 2, 1.

Trigonometric form of a complex number : VIII, 2, 2.

Trigonometric polynomials in \(m\) variables: \(\mathbf{X}\) , 4, 4.

Ultrametric space : IX, 2, Exercise 4.

Uniform approximation of a continuous real-valued function by continuous

functions belonging to a given set : X, 4, 1.

Uniform convergence in a subset of X : X, 1, 3.

Uniform convergence in the sets of \(\mathfrak{S}:\mathbf{X},\mathbf{1},\mathbf{2}\)

Uniform convergence, point of : X, 1, Exercise 9.

Uniform convergence, topology of: X, 1, 1.

Uniform convergence, uniformity of: X, 1, 1.

Uniformity, additive (of \(\mathbf{R}^n\) ): VI, 1, 2.

Uniformity of bounded convergence : X, 1, 3.

Uniformity of compact convergence : X, 1, 3.

Uniformity defined by a family of pseudometrics : IX, 1, 2.

Uniformity defined by a pseudometric : IX, 1, 2.

Uniformity of finite open coverings : IX, 4, Exercise 17.

Uniformity of pointwise convergence : X, 1, 3.

Uniformity of pointwise convergence in a subset of X : X, 1, 3.

Uniformity of precompact convergence : X, 1, 3.

Uniformity of \(\mathfrak{S}\) -convergence : X, 1, 2.

Uniformity of uniform convergence : X, I, I.

Uniformity of uniform convergence in a subset of X : X, 1, 3.

Uniformity of uniform convergence in the sets of S : X, 1, 2.

Uniformity, universal : IX, 1, Exercise 5.

Uniformly convergent (filter) : X, I, I.

Uniformly convergent in every set of \(\mathfrak{S}\) (family of mappings, filter, sequence, series): X, 1, 2.

Uniformly equicontinuous set of mappings : X, 2, 1.

Uniformly summable family : X, 1, 2.

Unit ball, sphere : VI, 2, 3 and IX, 3, 3.

Unit of angular measure : VIII, 2, 3.

Universal uniformity : IX, 1, Exercise 5.

Urysohn's theorems: IX, 4, 1 and 4, 2.

Valuation, real : IX, 3, 2.

Value, absolute : IX, 3, 2.

Value, absolute (of a complex number): VIII, 1, 1.

Value, semi-absolute : IX, 3, Exercise 9.

Valued division ring : IX, 3, 2.

Variety, coordinate : VI, 1, 4.

Vector, direction (of a line or ray): VI, 1, 4.

Vector subspaces, orthogonal : VI, 2, 2.

Vectors, orthogonal : VI, 2, 2.

Vertex of a stereographic projection : VI, 2, 4.

Weierstrass-Stone, theorem of : X, 4, 2.

Zero-dimensional space : IX, 6, 4
